\begin{scanpage}{0224}{206}
\noindent and finally
\[\begin{aligned}
&\frac{2\ln\dfrac{i}{\gamma ka}}{1+\psi(L)e^{ikL}}\bar\psi^2(L)e^{2ikL}\\
&=\left(1+i\frac32kL-k^2L^2-i\frac9{24}k^3L^3\right)\ln\frac{i}{\gamma ka}\\
&\quad-\frac{ikL}2\ln\frac{\gamma kL}i+\frac{ikL}2(1+\ln2)+\frac{k^2L^2}2\ln\frac{\gamma kL}i-\frac{k^2L^2}2(1+\ln2)\\
&\quad+i\frac5{24}k^3L^3\ln\frac{\gamma kL}i-ik^3L^3\left(\frac{14}{72}+\frac{5\ln2}{24}\right).
\end{aligned}\tag{33.16}\]
Using these relationships, it is not difficult to show that
\[F\left(\frac\pi2,\frac\pi2\right)=\frac{k^3L^3}{24}\left\{\frac1{\ln\dfrac La}+\frac{\dfrac73-\ln2}{\left(\ln\dfrac La\right)^2}+O\left[\left(\ln\frac La\right)^{-3}\right]\right\}.\tag{33.17}\]
The equation which has been found may be rewritten in the form
\[F\left(\frac\pi2,\frac\pi2\right)=\frac{k^3L^3}{24\left(\ln\dfrac{2L}a-\dfrac73\right)}\left\{1+O\left[\left(\ln\frac La\right)^{-2}\right]\right\}.\tag{33.18}\]
It completely agrees with Equation (33.06) which follows from [92].

This result confirms the correctness of scattering characteristic (32.05) calculated by us, and shows that it is applicable for vibrators of any length.
\sectiontitle{34}{The Results of Numerical Calculations}
The function $F(\vartheta,\vartheta_0)$ enables one to calculate the integral scattering thickness $S$ and the effective scattering area $\sigma$ of a passive vibrator. The integral scattering thickness is determined by the relationship
\[S=\frac Pp,\tag{34.01}\]
where
\[p=\frac{c}{8\pi}E_0^2\tag{34.02}\]
\end{scanpage}
